% Transcription — fonds Alexandre Grothendieck, shelfmark 114, batch 1 (pages 1-5).
% Established from the facsimile archives/batches/114/batch-01.pdf (a mirror of
% https://grothendieck.umontpellier.fr/114.pdf), per the skill
% transcribe-grothendieck.
% Five pages, all of them mathematics, and the whole folder in one batch.
% Pages 1 and 2 are a programme: nine numbered questions in English, each
% assigned in the margin to a chapter (IV to VII) and then triaged at the foot
% of the page into what will be treated, what will be deferred and what will be
% dropped; then a short list of the test categories in play. Page 3 is a
% typescript leaf of Pursuing Stacks, the author's page 350, carrying
% Proposition 4 on total W-asphericity — which is the statement pages 2, 4 and 5
% are circling. Pages 4 and 5, his own pages 1 and 2, are French notes on the
% aspheric part M_as of a modelizer, on when i* is conservative and faithful,
% and on two counterexamples. Page 2 is the verso of the typescript leaf.
% The transcription follows the language of each page, which is the author's
% own: English on 1, 2 and 3, French on 4 and 5.
% Pass: Opus 5 (claude-opus-5), 2026-08-17 — first pass, unchecked against the pages by a human.
\documentclass[11pt,a4paper]{article}
% Le préambule commun à toutes les transcriptions du fonds Grothendieck.
%
% Il tient en une idée : **l'incertitude se compose, elle ne se résout pas.**
% Une transcription qui lisse un mot douteux a détruit la seule chose pour
% laquelle on la faisait — un lecteur doit pouvoir savoir, à chaque endroit,
% ce qui était sur le feuillet et ce qui a été ajouté. Les six macros qui
% suivent sont donc l'appareil critique, et non de la décoration.
%
% Les mêmes commandes sont reconnues par `scripts/render.mjs`, qui produit la
% vue de lecture du site. Une macro ajoutée ici doit l'être aussi là-bas,
% faute de quoi le rendu échouera bruyamment — ce qui est voulu : un rendu
% silencieusement faux est pire qu'un rendu absent.

\ProvidesPackage{grothendieck}[2026/08/08 Transcriptions du fonds Grothendieck]

\usepackage{amsmath,amssymb,amsthm}
\usepackage{mathrsfs}
\usepackage{stmaryrd}
% Les diagrammes commutatifs sont partout dans ce fonds : c'est la langue
% même de la géométrie algébrique de Grothendieck, et une transcription qui
% les rendrait en prose ne transcrirait rien.
\usepackage{tikz-cd}
% Français par défaut — c'est la langue des feuillets — mais l'anglais est
% chargé aussi : la traduction et le résumé sont des documents anglais, et la
% typographie française (espaces avant les deux-points) y serait une faute.
% Ces documents basculent par \selectlanguage{english} en tête de corps.
\usepackage[english,french]{babel}
\usepackage{fontspec}
\usepackage{geometry}
\geometry{a4paper,margin=25mm,marginparwidth=28mm}
\usepackage{marginnote}
\usepackage{xcolor}
% ulem, not soul. `soul`'s \st and \ul are built on a character-by-character
% scan that XeLaTeX's font machinery breaks: the document fails with an
% undefined control sequence rather than degrading. ulem does the same two jobs
% and is Unicode-engine safe. `normalem` stops it hijacking \emph, which the
% transcriptions use for Grothendieck's own emphasis.
\usepackage[normalem]{ulem}
\usepackage{hyperref}
\hypersetup{colorlinks=true,linkcolor=black,urlcolor=black,pdfborder={0 0 0}}

% --- Métadonnées du lot -----------------------------------------------------
% Reprises telles quelles par le script de rendu, qui les affiche en tête de
% la vue de lecture. Elles fixent ce que le fichier prétend couvrir : sans
% elles, une transcription flotte sans point d'attache dans un fonds de seize
% mille feuillets.

\newcommand{\folder}[1]{\gdef\grothFolder{#1}}
\newcommand{\batch}[1]{\gdef\grothBatch{#1}}
\newcommand{\leaves}[2]{\gdef\grothFirst{#1}\gdef\grothLast{#2}}
\newcommand{\foldertitle}[1]{\gdef\grothFoldertitle{#1}}
\gdef\grothFolder{?}\gdef\grothBatch{?}\gdef\grothFirst{?}\gdef\grothLast{?}\gdef\grothFoldertitle{}

% --- Le résumé --------------------------------------------------------------
%
% Il ouvre la lecture modernisée, sous le titre « Résumé », et il est composé
% en petit. Ce n'est pas un ornement : c'est le seul passage du document qui
% ne soit pas des mathématiques, et un lecteur doit pouvoir le reconnaître
% avant de l'avoir lu — puis le sauter s'il connaît déjà le sujet, ou s'y
% arrêter s'il ne le connaît pas. Le filet qui le suit marque la reprise du
% corps.

\newenvironment{resume}
  {\small\setlength{\parskip}{0.45em}}
  {\par\medskip\hrule\medskip}

% --- Mots-clés ---------------------------------------------------------------
%
% En anglais, en clôture du résumé de la lecture modernisée : le vocabulaire
% moderne sous lequel chercher ce que les feuillets construisent. Le manifeste
% du site (`npm run manifest`) les extrait pour étiqueter la cote entière —
% cette ligne est la source unique des tags, il n'y a pas de fichier à côté.
\newcommand{\keywords}[1]{\par\smallskip\noindent
  {\footnotesize\textsc{Keywords} --- \textit{#1}}\par}

% --- L'appareil critique ----------------------------------------------------

% Le numéro de feuillet, en marge. C'est l'ancre qui permet de régler un
% désaccord sur une formule en regardant *un* feuillet plutôt qu'une chemise
% de six cents. Le site s'en sert aussi pour tourner les pages du fac-similé
% quand on fait défiler la transcription.
\newcommand{\leaf}[1]{%
  \marginnote{\footnotesize\color{gray!70}\textsf{#1}}%
  \label{leaf:#1}\ignorespaces}

% Illisible : on ne devine pas. Un mot inventé qui se lit comme les autres est
% le pire résultat possible.
\newcommand{\ill}{\textcolor{red!60!black}{[\,\ldots\,]}}

% Lecture douteuse : le mot est donné, et signalé comme incertain.
\newcommand{\uncertain}[1]{\uline{#1}}

% Ajout de l'éditeur — une lettre manquante, un mot sous-entendu.
\newcommand{\add}[1]{\textcolor{blue!50!black}{[#1]}}

% Biffé par Grothendieck lui-même. Ce qu'il a rayé fait partie du document :
% une rature dit souvent où la pensée a changé de direction.
\newcommand{\struck}[1]{\sout{#1}}

% Note du transcripteur, sur l'état matériel du feuillet.
\newcommand{\note}[1]{\par\smallskip{\small\color{gray!60!black}\textit{[#1]}}\par\smallskip}

% Note marginale de Grothendieck, distincte du corps du texte.
\newcommand{\marginal}[1]{\par\smallskip\noindent\hspace{1em}%
  {\small\color{gray!60!black}#1}\par\smallskip}

% --- Titre ------------------------------------------------------------------

\AtBeginDocument{%
  \begin{center}
    {\large\bfseries Fonds Alexandre Grothendieck --- cote n\textsuperscript{o}~\grothFolder}\\[2pt]
    {\itshape\grothFoldertitle}\\[4pt]
    {\small feuillets \grothFirst--\grothLast\ (lot \grothBatch)}\\[2pt]
    {\footnotesize\color{gray!60!black}Transcription établie d'après le fac-similé numérisé
     par l'Université de Montpellier.\\
     Les lectures incertaines sont soulignées, les illisibles notées [\,\ldots\,],
     les ajouts entre crochets.}
  \end{center}
  \vspace{1em}}

\endinput


\folder{114}
\batch{1}
\pages{1}{5}
\dating{[1983]}
\watermark{Édition de démonstration}
\foldertitle{Asphéricité : notes manuscrites (s.d.)}

\begin{document}

\page{1}

En haut à droite, d'une encre plus pâle, le titre de la chemise :
« Asphéricité ».

\begin{enumerate}
\item[1)] Morphisms of asph.\ structures ? \ill\ images ?

\item[2)] Contr\add{a}ctors, morphisms of such

\item[3)] Induced structures on \uncertain{$M_{/a}$}

\item[4)] Lack of autoduality

\item[5)] When does \struck{an} \uncertain{lg.}\ structure type give rise to a
can.\ modelizer ?

\item[6)] Existence theorems for test functors

\item[7)] Finiteness conditions in an elementary
\end{enumerate}

\note{le point 7 s'arrête sur ce mot ; le substantif attendu n'est pas écrit.
Le numéro 7 est tracé par-dessus un 6 : les deux points 6 et 7 sont d'abord
venus sous le même numéro, et une accolade les réunit dans la marge de droite}

\begin{enumerate}
\item[8)] Computation of \uncertain{cohomology} by boundary operator ?

\item[9)] $\Delta^{f}$, infinite \ill\ and finite \ill\
\uncertain{semi-}modelizers of (Cat).
\end{enumerate}

\marginal{VI}

\marginal{V ? Also VII}

\note{ces deux mentions sont portées dans la marge de droite, la première en
regard du point 3, la seconde en regard des points 5 et 6 ; ce sont des numéros
de chapitre, du même ordre que le « part IV » du bloc qui suit}

Au bas de la page, entre parenthèses :

\begin{quote}
(\struck{7} 7) and 8) may not be treated at all,

2, 3, 9 should be done in part IV

1, 4, 5, 6 can be looked at Caen.)
\end{quote}

\note{les trois lignes se partagent exactement les neuf points de la liste :
7 et 8 abandonnés, 2, 3 et 9 renvoyés à la partie IV, 1, 4, 5 et 6 remis à
Caen. C'est cette partition complète qui rend la lecture des chiffres sûre,
là où le 3 de la deuxième ligne et le 5 de la troisième se distinguent mal
l'un de l'autre}

\page{2}

$\Delta^{f}$

$\square$, $\square^{f}$

$\bigcirc$, $\bigcirc^{f}$

$\mathcal{B}_{f}(I)$ weak test category ?

\marginal{(Ord) has asph.\ str.\ or contr.\ str.}

cf.\ p.\ 74

total 0-connectedness $\Longrightarrow$ total asphericity ?

\note{le trait de l'implication ne porte de pointe qu'à droite : c'est une
implication simple et non une équivalence. La question de cette dernière ligne
est exactement celle que tranche la proposition 4 de la page suivante}

\page{3}

\note{feuillet dactylographié, paginé 350 de la main de l'auteur : une page de
\emph{Pursuing Stacks}. Les passages biffés le sont à la machine, par
surfrappe de x ; les ajouts en interligne sont manuscrits. La page est
transcrite ici parce qu'elle porte la proposition dont les pages manuscrites
de la chemise discutent}

(iii) For any $x$ in $C$, $i^{*}(x)$ is 0-connected in $\hat{A}$ i.e.\
\[
i^{*}(C) \ \ill\ \ \hat{A}_{W_{0}} \quad
(=_{\mathrm{def}} \text{set of 0-connected objects in } \hat{A}).
\]

\note{le signe d'inclusion entre $i^{*}(C)$ et $\hat{A}_{W_{0}}$ manque : la
machine a laissé l'espace et rien n'y a été porté à la main}

Moreover, when these conditions are satisfied, \struck{and if x x then} and if
the contractibility structure of $M$ is non-trivial \struck{iff the one of
$\hat{A}$ is} \struck{and x} then so is the one of $A$ and $A$ is a
\add{strict} test category.

Proof : \struck{the last} It remains to prove the last statement, which follows
from corollary 1. Here by « strict test category » we mean strict $W$-test
category, which implies that the contractor $A$ is non trivial visibly. But of
course, as condition (i) does not depend on $W$, we may take in the last
conclusion $W = W_{\infty}$, which is the strongest --- but we got this already
by prop.\ 3 directly, for any non trivial contractor. Thus the complement about
$A$ being a test or strict test category is interesting mainly, when we do not
assume beforehand that $A$ is a contractor. So I'll have to clarify the points
a) and b) of remark 2 above.

After a minute\add{'}s thought, I feel very silly to find out that point a) is
about trivial --- that I have been dragging for months the notion of a
total\add{ly} $W$-aspheric category $A$, without taking the trouble to sit down
and look whether it really depends on $W$, and at the same time, whether the
condition that the products $a \times b$ in $\hat{A}$ (for $a$ and $b$ in $A$)
are $W$-aspheric, already implies that $A$ is $W$-aspheric and hence totally
$W$-aspheric ! So let's clean it up in the long last !

\textbf{Proposition 4.} Let $A$ be any small category, $W$ a basic localizer
satisfying Loc 4), $W_{\infty}$ the finest basic localizer (i.e.\ usual weak
equivalence). The following conditions are equivalent :

\begin{enumerate}
\item[(i)] $A$ is totally $W_{\infty}$-aspheric.

\item[(ii)] $A$ is totally $W$-aspheric (cf.\ section 65 A), def.\ 1 p.\ 170)).

\item[(iii)] $A$ is non empty, and for any two objects $a$, $b$ in $A$, the
product $a \times b$ in $\hat{A}$ is 0-connected.
\end{enumerate}

Proof. From the definition (loc.\ cit.) it is clear that
\[
\text{(i)} \quad \ill \quad \text{(ii)} \quad \ill \quad \text{(iii)},
\]
and it is immediate that (iii) implies that $A$ is 0-connected, hence (iii)
also means that $A$ is totally $W_{0}$-aspheric, where $W_{0}$ is the coarsest
basic localizer satisfying Loc 4), namely

\note{les signes d'implication de cette ligne manquent eux aussi : la machine a
ménagé les intervalles et ils sont restés vides. La phrase s'interrompt sur
« namely » — la suite est à la page 351, qui n'est pas dans la chemise}

\page{4}

\begin{tikzcd}[column sep=large]
\hat{A} & M \arrow[l, "i^{*}"'] & M' \arrow[l, "f"']\\
 & M_{as} & M'_{as}
\end{tikzcd}

\note{les deux mentions de la seconde ligne sont portées sous $M$ et sous $M'$,
sans flèche : ce sont des annotations et non des objets d'un diagramme}

\begin{tikzcd}
A \arrow[r, "i"] \arrow[dr, dashed] & \underline{M}\\
 & M_{as} \arrow[u]
\end{tikzcd}

$i$ $M_{as}$-asphérique $\overset{?}{\Longrightarrow}$ $i^{*}$ conservatif et
fidèle

\underline{NB.} Si $A \xrightarrow{f} B$ asphérique, alors
$f^{*} : \hat{B} \to \hat{A}$ conservatif et fidèle, i.e.\ $f$ dominant, i.e.\
le cocrible engendré par \ill\ $f^{*}$ \ill\ propriétés d'exactitude de
$f^{*}$ \ill

\marginal{i.e.\ $\forall\, b \in \mathrm{Ob}\, B$, $\exists\, a \in A$ et
$f(a) \to b$, i.e.\ $A_{/b} \neq \emptyset$. OK}

\struck{que \ill\ $b \to b'$ tel que \ill\ $f^{*}(b) \to f^{*}(b')$ \ill}

\[
\bigl[\, f^{*}(x) = \emptyset_{\hat{A}} \Longrightarrow x =
\emptyset_{\hat{B}} \,\bigr]
\]

\struck{$A_{/f(b)}$} \quad $A_{/f(x)}$

Mais $f^{*}(a) = \emptyset_{\hat{B}} \Longrightarrow
\mathrm{Hom}_{\hat{B}}(b, a) = \emptyset$ pour tout \ill\
$b \in f(\mathrm{Ob}(A))$,

\struck{Il faut en déduire que $\mathrm{Hom}(b,a) = \emptyset$ $\forall\, b
\in B$}

\struck{Plus \ill} \add{ou conservatif}

Pour conclure que $i^{*}$ est fidèle, il faut savoir que
$M \to \hat{M_{as}}$ (i.e.), il faut que $M_{as}$ est \underline{génératrice}.

Peut-on avoir $M_{as} = \emptyset$ / 1) \ill\ que $A = \emptyset$, donc
$\hat{A} = \{e\}$, donc $\hat{A}_{as} = \emptyset$, \struck{ok}

Contre-exemple ! Peut-on avoir $M_{as} = \{a\}$

Alors \uncertain{DPS} $A = \{a\}$, $\hat{A} = \ill$

\struck{$M_{as} = \{x \in \mathrm{Ob}\,M \mid \mathrm{Hom}\,(a,x) \ill\}$}

\note{la dernière ligne de la page, ajoutée sous celle-ci, est coupée par le
bord inférieur du feuillet et n'a pas été rétablie ; elle reprend au haut de la
page suivante}

\page{5}

\note{le feuillet porte « 2 » en haut à droite : c'est la pagination de
l'auteur, la page précédente étant sa page 1}

Soit $G = \underline{\mathrm{End}}(a)$, $\hat{A} \simeq (G\text{-ens})$
\uncertain{donc nécessairement} $G = \{\mathrm{id}\}$, alors
$\hat{A} \simeq (\mathrm{Ens})$, $\hat{A}_{as} = \{$\uncertain{ens.\
uniponctuels}$\}$,

donc $M_{as} = \{x \in \mathrm{Ob}\,M \mid \mathrm{Hom}\,(b,x)$ card 1$\}$

\struck{donc} Contre-exemple : $M$ cat.\ discrète de cardinal $\geq 2$, !

$M_{as}$ engendre $M$ de \uncertain{façon} que $M \to \hat{M}_{as}$ soit
fidèle et conservatif.

\begin{tikzcd}[column sep=large]
\hat{A} \arrow[r, bend left=20] & M \arrow[l, bend left=20] & M' \arrow[l, "f^{*}"']
\end{tikzcd}

$f^{*}$ commute aux limites

$f_{!}$ existe alors, si \ill\ $\underline{W}$ \ill\ que tout foncteur sur $M'$
qui commute aux limites soit représentable --- à \uncertain{vérifier} des
conditions qui l'assurent, qu'il \ill\ pas \ill\ : faire des \uncertain{hyp.}\
d'\ill\ \uncertain{accessibilité} sur $f^{*}$. Mais \ill\ que $f_{!}$ \ill\
$M_{as}$ dans $M'_{as}$ ?

\begin{tikzcd}[column sep=large]
M \arrow[r, bend left=20, "f_{!}"] & M' \arrow[l, bend left=20, "f^{*}"]
\end{tikzcd}

\note{la question de la dernière ligne — si l'adjoint à gauche $f_{!}$ respecte
les parties asphériques — est celle vers laquelle toute la page converge, et
elle reste ouverte : rien ne suit sur le feuillet}

\end{document}
